% Created 2026-04-15 Wed 17:52
% Intended LaTeX compiler: pdflatex
\documentclass[11pt]{article}

\usepackage[utf8]{inputenc}
\usepackage[T1]{fontenc}
\usepackage{graphicx}
\usepackage{longtable}
\usepackage{wrapfig}
\usepackage{rotating}
\usepackage[normalem]{ulem}
\usepackage{amsmath}
\usepackage{amssymb}
\usepackage{capt-of}
\usepackage{hyperref}
\author{Patrick Clark}
\date{\today}
\title{Intro}
\hypersetup{
 pdfauthor={Patrick Clark},
 pdftitle={Intro},
 pdfkeywords={},
 pdfsubject={},
 pdfcreator={},
 pdflang={English}}
\begin{document}

\maketitle
\tableofcontents

\section{Alignment\hfill{}\textsc{ATTACH}}
\label{sec:org43717f6}
Allineare le sequenze è importante per ottenere una comparazione diretta, oltre ad essere prerequisito per analisi più complesse contro database or per la costruzione di alberi filogenetici o identificazione di motivi funzionali
\textbf{Omologia} -> Due entità condividono la stessa origina filogenetica, da cui si sono evoluti, differenziandosi dalle altre -> \textbf{Qualitativa}
\textbf{Similarità} -> Il termine similarità ha un significato generale, utilizzato a prescindere dalla motivazione. La similarità biologica può essere dovuta ad omologia ma anche dovuta al caso, o da fenomeni di convergenza adattiva -> \textbf{Quantitative}
\begin{itemize}
\item Ortologhi -> Speciazione
\item Parologhi -> Duplicazione
\item Xenologhi -> trasferimento orizzontale
\textbf{\textbf{BioMart}} -> Raccoglimento di ortologhi e parologhi di un gene
\end{itemize}
\begin{center}
\includegraphics[width=.9\linewidth]{/home/patrick/org/.attach/7f/3e9447-83b3-488d-a0b9-d68f3a5d5a42/Screenshot from 2026-04-14 17-11-14.png}
\end{center}
\textbf{Gap Opening} -> Si apre un gap assegnando uno score negativo, la cui estensione comporta ulteriore penalità
\subsection{Matrici}
\label{sec:org3e29e39}
\begin{enumerate}
\item Identità
\item Genetic-Code
\item Matrice delle proprietà chimiche
\item Matrici Empiriche
\begin{enumerate}
\item PAM
\begin{enumerate}
\item BLOSUM
\begin{enumerate}
\item Gonnet
\end{enumerate}
\end{enumerate}
\end{enumerate}
\end{enumerate}
PAM e BLOSUM si basano sulle frequenze osservate di specifiche sostituzioni amminoacidiche. La similarità tra due amminoacidi è proporzionale alla loro frequenza di sostituzione osservata.
Queste matrici si basano su evidenze biologiche, le cui differenze osservata tra sequenze omologhe negli allineamenti è dovuta solo a eventi di mutazione. Alcune di queste mutazioni hanno effetti minimi o assenti nella struttura/funzione delle proteine.

La scelta di una delle matrici è cruciale per ottenere buoni risultati.
Per sequenze meno divergenti, BLOSUM80 e PAM1 sono le migliori. BLOSUM62 e PAM120 per mediamente divergenti. BLOSUM45 e PAM250 per sequenze molto divergenti.
\subsubsection{PAM}
\label{sec:org197de09}
L'idea base delle matrici PAM nasce da studi di filogenetica molecolare su 71 famiglie di proteine con un minimo di 85\% di similarità. L'assunzione di partenza è che analizzando sequenze filogeneticamente associate è possibile calcolare la percentuale di mutazioni accettate. Una unità di PAM è definita come l'1\% di cambiamento di posizione di un amminoacido, cioè una sostituzione ogni 100 amminoacidi.
Si è partiti ricostruendo la sequenza ancestrale partendo da allineamento multiplo tra proteine con similarità elevata >85\%.
Le sostituzioni ``accettate'' sono riportate su una matrice da cui è possibile ottenere la mutabilità relativa (\(m_i\)) del residuo i.
Per ogni amminoacido è quindi possibile valutare la possibilità di cambiamento.
Una volta definita la mutabilità, la domanda da porsi è qual è la probabilità che sia cambiato in j.
$$M_{ij} = m_i \cdot \frac{f_{ij}}{f_i} = \frac{f_{ij}}{f_i} \cdot \frac{f_i}{100 \cdot fp_i} = \frac{f_{ij}}{100 \cdot fp_i}$$
La probabilità che il residui i non cambi è \$1-M\textsubscript{ij}\$\textsubscript{nil}.
La probabilità di una mutazione dipende dal tempo trascorso dalla divergenza delle sequenze. I valori sono normalizzati per il tempo necessario ad introdurre 1 mutazione ogni 100 residui (distanza = 1PAM).
Matrici PAM con indice maggiore di 1 si ottengono moltiplicando la matrice PAM per se stessa k volte, PAM2 = PAM1 x PAM1
Le Matrici PAM si dividono in
\begin{itemize}
\item Scoring
\item Substitution
\end{itemize}
Le matrici di scoring utilizzate dai programmi di allineamento sono calcolate dalle matrici di sostituzione come:
$$s(a,b) = int(10\cdot log{\frac{M(a,b)} {C(a,b)}})$$
Dove
\begin{itemize}
\item s(a,b) -> score per una coppia di amminoacidi
\item M(a,b) -> probabilità di sostituzione di un amminoacido
\item C(a,b) -> probabilità di una sostituzione di un amminoacido
\end{itemize}
Il rapporto \(\frac{M(a,b)}{C(a,b)}\)  indica quante volte la probabilità di omologia è maggiore della probabilità casuale.
\subsubsection{BLOSUM}
\label{sec:org34d20a9}
Le matrici BLOSUM si basano sul BLOCKS database, il quale mantiene blocchi di sequenze relazionate tra loro con un definito valore di identità amminoacidica, generalmente tra il 30\% ed il 95\%. A partire da ogni blocco è possibile calcolare la frequenza di sostituzione osservata e la probabilità di accoppiamento casuale
$$s(a,b) = int(10\cdot log{\frac{M(a,b)} {C(a,b)}})$$
In questo caso \(M(a,b)\) si riferisce alla probabilità di sostituzione osservata per ogni blocco.
\begin{itemize}
\item Per ogni allineamento nel database BLOCKS le sequenze sono raggruppate in cluster con almeno il 50\% (per BLOSUM50, altrimenti 80\% per BLOSUM80 ecc\ldots{}) di residui identici
\item Tutte le coppie di sequenze sono comparate, e la frequenza delle coppie osservata vengono annotate
\item La frequenza attesa delle singole coppie vengono computate dalle singole frequenze amminoacidiche $$f_{A,C} = f_A \cdot f_C$$
\begin{itemize}
\item Per ogni coppia di amminoacidi lo score di sostituzione viene computato essenzialmente con \(s(a,b)\).
\end{itemize}
\end{itemize}
\subsection{DotPlot\hfill{}\textsc{ATTACH}}
\label{sec:org190018b}
Il dotplot aiuta a ispezionare la relazione tra due sequenze. Aiuta ad identificare potenzialmente regioni conservate. Si ottiene in 4 maniere:
\begin{itemize}
\item Diagonalmente attraverso i dot -> Match
\item Diagonalmente attraverso celle vuote -> Mismatch
\item Orizzontalmente -> Gap
\item Verticalmente -> Gap
\end{itemize}
L'allineamento viene effettuato muovendosi lungo la matrice maniera tale da utilizzare le diagonali con i punti più grandi. Questo massimizza il numero di match.
Il metodo DotPlot si utilizza sia su sequenze nucleotidiche sia su sequenze proteiche.
Poichè all'aumentare della lunghezza della sequenza aumentano anche i match spuri e quindi anche il rumore di fondo, i parametri del window size e della stringency aiutano a visualizzare meglio.
\begin{center}
\includegraphics[width=.9\linewidth]{/home/patrick/org/.attach/7f/3e9447-83b3-488d-a0b9-d68f3a5d5a42/Screenshot from 2026-04-15 16-21-31.png}
\end{center}
\subsection{Tipologie di Allineamento}
\label{sec:org255f089}
\begin{itemize}
\item \textbf{Globale} -> L'allineamento globale tiene conto della similarità tra sequenze per l'intera lunghezza (Query Coverage 100\% sempre)
\item \textbf{Locale} -> L'allineamento locale tiene conto della similarità tra sequenze soltanto per specifiche regioni simile tra parti della sequenza in analisi (Cambia il Query Coverage)
\end{itemize}
Importante considerare che \textbf{lo score di allineamento più elevato non significa necessariamente che sia biologicamente significativo rispetto ad un altro allineamento con score inferiore}.
Il tipo di allineamento (globale o locale) dipende strettamente dal caso in studio. A volte è meglio un allineamento locale, altre volte uno globale.
Considerando che il numero di allineamenti cresce al crescere della lunghezza delle sequenze, tecniche di programmazione come il dynamic programming vengono utilizzate per esplorare le possibili soluzioni al problema in un tempo ragionevole.
L'algoritmo globale \textbf{Needleman-Wunsch} e quello locale \textbf{Smith-Waterman} utilizzano questa tecnica.
\subsubsection{Needleman-Wunsch}
\label{sec:org9f1865a}
L'algoritmo di Needleman-Wunsch può essere riassunto in 3 step principali:
\begin{enumerate}
\item Riempimento della matrice
\item Calcolo dello score
\item Allineamento
\end{enumerate}
Le sequenze compongono la matrice principale e le regole di composizione sono uguali a quelle del DotPlot.
Ogni punto della matrice può essere raggiunto da 3 possibili posizioni. Si sceglie quello con lo score più elevato tra le 3 possibilità per ottenere il best alignment score per uno specifico punto.
\subsubsection{Smith -Waterman}
\label{sec:org2b76f64}
La Smith-Waterman la si può ottenere modificando la formula introducendo lo 0 cosiì da non avere score negativi.
L'allineamento si ricostruisce dal valore più elevato e l'allineamento termina allo 0.
Lo score sarà infatti pari al massimo tra score(x,y-1) - gap\textsubscript{penalty}, score(x-1, y-1) + substitution-score(x,y), score (x-1, y) - gap-penalty
In principio l'allineamento locale è utile per cercare similarità in un databse. Smith-Waterman è troppo lento però per identificare regioni simili ad una query in database molto grandi. L'algoritmo Smith-Waterman genera allineamenti ottimali ed esaustivi unicamente.
Per la ricerca su database si utilizzano algoritmi più complessi quali FASTA e BLAST.
\subsection{FASTA}
\label{sec:org36a67ff}
L'algoritmo FASTA viene sviluppato nel 1985 basandosi sull'indicizzazione delle parole di una data lunghezza (k-tuple) (2 per a.a. e 6 per acidi nucleici). L'indicizzazione consiste nel creare una lista di tutte le posizioni in cui ogni possibile w (word) appare all'interno della sequenza. Nel caso delle proteine, per k =2 avremo 40 parole diverse (\(20\cdot N\) dove N è la lunghezza delle k-tuple).
Il primo step dell'algoritmo FASTA è quello di creare un indice di word (w) della sequenza di query. Ogni parola verrà poi utilizzata per la ricerca all'interno del database
Nello specifico:
\begin{itemize}
\item Troviamo parole contigue sulla stessa diagonale
\item Calcoliamo gli score con le matrici di sostituzione
\item Si uniscono i frammenti che possono essere uniti con una threshold T di accettabilità
\item Si applica Smith-Waterman su una porzione solamente (i frammenti riuniti) per ottenere l'allineamento migliore
\end{itemize}
\subsection{BLAST}
\label{sec:orge80e566}
\begin{enumerate}
\item Si crea quindi un indice di tutte le w di una sequenza di query (con overlap)
\begin{itemize}
\item Si ottengono tutte le w dalla sequenza di query
\item Si va a creare una lista di tutte le query che hanno un match al di sopra della threshold T (se la w è GSW -> GTW GNW GAW ATW GTF sono alcune delle possibili word. In base alla sostituzione si calcola lo score e se nonostante la sostituzione siamo sopra T, allora viene aggiunta la word alla lista)
\end{itemize}
\item Ogni parola dalla lista compilata (quindi incluse quelle generate in base alle matrici di sostituzione) viene ricercata all'interno del database per trovare i match all'interno del DB
\item Si estende la sequenza in ambe due direzioni
\item Ci si ferma quando lo score diminuisce, o la sequenza termina (full-match)
\end{enumerate}
\end{document}
